\chapter{Discussion}
\label{sec:discussion}

\section{When to use PACE}
\label{sec:d-when}

The results give a clear rule for when PACE is the right tool. It is preferred on
systems with broad priors, high information per experiment, and stiff or
non-differentiable dynamics, in low to moderate-dimensional continuous design
spaces. These are exactly the systems on which the prior-contrastive training
signal saturates, and the screening threshold $L_{\mathrm{escape},50}$ identifies
them before any policy is trained. On the saturated anchors, CES, SEIR R=3, and
ASM1, PACE matches or beats the policy, and on ASM1 it works where policy training
is infeasible.

Learned policies remain the better choice elsewhere. When the escape threshold is
small, the policy trains without saturation and its trajectory-level view can
capture design dependencies that a greedy per-step rule does not. This is what we
see on Monod, on PI3K, and on the highest-dimensional source location task. The
crossover is governed by saturation, not by parameter dimension and not by the
effective sample size, so the two methods are complementary rather than competing.

\section{Practical cost}
\label{sec:d-cost}

PACE shifts cost from a per-step policy-training budget to a one-time inference
pass plus a modest per-step search. Training the neural posterior estimator takes
between about $3.6$ and $12.8$ minutes once per system. At deployment, PACE selects
each design in under four seconds on eleven of the twelve systems, the exception
being the stiff ASM1 model at about forty seconds per step. By contrast, training a
policy ranges from a few minutes on the simplest systems to about sixteen hours on
Monod, and is unstable on CES and infeasible on ASM1 and on SEIR R=3 at the budget
needed to escape saturation. A deployed policy then produces designs in
milliseconds, which is faster than PACE per step, but this difference is irrelevant
for scientific experiments that take minutes to weeks to run. The full timing table
is in Appendix~\ref{app:walltime}.

\section{Limitations}
\label{sec:d-limits}

Several limitations qualify the conclusions. The forward-only search does not scale
to high-dimensional or combinatorial design spaces, where a learned policy is
stronger; the source location result at $p{=}8$ is the clearest case, and the
ablation shows the optimiser rather than the estimator is the bottleneck there.
Under a collapsed effective sample size, PACE produces reliable design rankings but
not calibrated EIG values, so it should be read as an acquisition criterion rather
than an absolute information estimate. The method assumes the neural posterior
estimator does not collapse onto the wrong mode; this is benign for the unimodal
posteriors of well-observed ODE systems but is a genuine risk under strong
multimodality, and it must be checked with the diagnostics of
Section~\ref{sec:r-diag}. The Gaussian-linearised threshold $\Lambda$ is a
qualitative guide, not a quantitative predictor, on stiff systems. Finally, the
study evaluates greedy design and a fixed myopic horizon; non-myopic extensions are
left open.

On reproducibility and validity, all main comparisons use $n=300$ paired rollouts
with paired tests and report both bounds on the EIG, so the orderings marked
definitive cannot be reversed by estimation slack. Earlier small-sample estimates
were replaced by these larger paired runs.

\section{Broader impact}
\label{sec:d-impact}

More efficient experimental design reduces the number of experiments needed to
learn a system, which lowers the cost, time, and material or animal use of
scientific work in areas such as pharmacokinetics, wastewater treatment, and
bioprocessing. The same machinery applies to epidemic intervention design, where
recommendations can affect public health policy and should be used with appropriate
caution and human oversight. PACE also makes its own uncertainty visible through the
diagnostics it reports, which supports honest use under model misspecification
rather than overconfident automation.
